% !TEX program = pdflatex
\documentclass[letterpaper,10pt]{article}

\usepackage{latexsym}
\usepackage[empty]{fullpage}
\usepackage{titlesec}
\usepackage{marvosym}
\usepackage[usenames,dvipsnames]{color}
\usepackage{verbatim}
\usepackage{enumitem}
\usepackage[hidelinks]{hyperref}
\usepackage{fancyhdr}
\usepackage[english]{babel}
\usepackage{tabularx}
\input{glyphtounicode}

\pagestyle{fancy}
\fancyhf{}
\fancyfoot{}
\renewcommand{\headrulewidth}{0pt}
\renewcommand{\footrulewidth}{0pt}

\addtolength{\oddsidemargin}{-0.5in}
\addtolength{\evensidemargin}{-0.5in}
\addtolength{\textwidth}{1in}
\addtolength{\topmargin}{-.5in}
\addtolength{\textheight}{1.0in}

\urlstyle{same}
\raggedbottom
\raggedright
\setlength{\tabcolsep}{0in}

\titleformat{\section}{
  \vspace{-6pt}\scshape\raggedright\large
}{}{0em}{}[\color{black}\titlerule \vspace{-5pt}]

\pdfgentounicode=1

\newcommand{\resumeItem}[1]{
  \item\small{#1 \vspace{-2pt}}
}

\newcommand{\resumeSubheading}[4]{
  \vspace{-2pt}\item
    \begin{tabular*}{0.97\textwidth}[t]{l@{\extracolsep{\fill}}r}
      \textbf{#1} & #2 \\
      \textit{\small#3} & \textit{\small #4} \\
    \end{tabular*}\vspace{-7pt}
}

\newcommand{\resumeProjectHeading}[2]{
    \item
    \begin{tabular*}{0.97\textwidth}{l@{\extracolsep{\fill}}r}
      \small#1 & #2 \\
    \end{tabular*}\vspace{-7pt}
}

\newcommand{\resumeSubItem}[1]{\resumeItem{#1}\vspace{-4pt}}
\renewcommand\labelitemii{$\vcenter{\hbox{\tiny$\bullet$}}$}
\newcommand{\resumeSubHeadingListStart}{\begin{itemize}[leftmargin=0.15in, label={}]}
\newcommand{\resumeSubHeadingListEnd}{\end{itemize}}
\newcommand{\resumeItemListStart}{\begin{itemize}}
\newcommand{\resumeItemListEnd}{\end{itemize}\vspace{-5pt}}

\begin{document}

\begin{center}
    \textbf{\Huge \scshape Nico Ohayon} \\ \vspace{1pt}
    \small 469-560-0141 $|$ \href{mailto:nohayonr@gmail.com}{\underline{nohayonr@gmail.com}} $|$
    \href{https://linkedin.com/in/nico-ohayon-rozanes}{\underline{linkedin.com/in/nico-ohayon-rozanes}} $|$
    \href{https://github.com/NicoOhR}{\underline{github.com/NicoOhR}} $|$
    \href{https://nicoohr.github.io/Sicarii}{\underline{Portfolio}}
\end{center}

\section{Education}
  \resumeSubHeadingListStart
    \resumeSubheading
      {University of Texas at Dallas}{Richardson, TX}
      {Bachelor of Science in Data Science, Double Degree in Mathematics}{Aug. 2023 -- May 2027}
  \resumeSubHeadingListEnd

\section{Experience}
  \resumeSubHeadingListStart

    \resumeSubheading
      {Software Engineering Intern}{May 2025 -- Aug 2025}
      {Exolambda}{Richardson, TX}
      \resumeItemListStart
        \resumeItem{Built a multi-layer sensor telemetry pipeline (ADC/IMU) and communication stack from scratch, reducing end-to-end latency by 40\% and increasing throughput.}
        \resumeItem{Standardized the embedded build system and deployment workflow, achieving reproducible builds and reducing provisioning complexity.}
        \resumeItem{Collaborated with hardware, architecture, and controls teams to validate system integration and flight algorithms; cleared firmware blockers ahead of delivery deadlines.}
      \resumeItemListEnd

    \resumeSubheading
      {Bioinformatics Research Assistant}{August 2025 -- Present}
      {University of Texas at Dallas}{Richardson, TX}
      \resumeItemListStart
        \resumeItem{Built a high-throughput ML data processing pipeline over large-scale mouse IMU datasets; accelerated pairwise similarity computation with CUDA from 12 hours to 3 hours.}
        \resumeItem{Designed reproducible experiment tracking infrastructure for systematic parameter sweeps with auditable, version-controlled results.}
        \resumeItem{Eliminated correctness failures in a legacy clustering pipeline and developed time-series preprocessing that measurably improved model output interpretability.}
      \resumeItemListEnd

    \resumeSubheading
      {Embedded Subteam Lead}{December 2023 -- December 2024}
      {Dallas Formula Racing}{Richardson, TX}
      \resumeItemListStart
        \resumeItem{Led firmware and electronics for a 16-sensor real-time data acquisition and telemetry system; ran weekly design reviews for 12 engineers through integration and track testing.}
      \resumeItemListEnd

    \resumeSubheading
      {Undergraduate Research Scholar Award}{August 2024 -- December 2024}
      {University of Texas at Dallas}{Richardson, TX}
      \resumeItemListStart
        \resumeItem{Developed high-performance C++ combinatorial search algorithms with state-space pruning for large-scale automated experiments; analyzed outputs to validate results and surface new conjectures.}
      \resumeItemListEnd

  \resumeSubHeadingListEnd

\section{Projects}
    \resumeSubHeadingListStart

      \resumeProjectHeading
          {\textbf{FSAEStats} $|$ \emph{Rust, DuckDB, Apache Arrow}}{}
          \resumeItemListStart
            \resumeItem{Built an automated data engineering pipeline extracting structured data from raw PDFs into a queryable backend; designed a high-performance Rust analytical layer using Arrow and DuckDB for fast aggregations over large datasets.}
          \resumeItemListEnd

      \resumeProjectHeading
          {\textbf{Three Body Simulation} $|$ \emph{Rust, Tokio, Python, gRPC, Bevy}}{}
          \resumeItemListStart
            \resumeItem{Built a distributed simulation system in Rust with a gRPC interface enabling Python clients to orchestrate and automate large experimental campaigns across the compute layer.}
          \resumeItemListEnd

      \resumeProjectHeading
          {\textbf{TradeX} $|$ \emph{Python, PyTorch, NumPy, Docker}}{January 2024 -- May 2024}
          \resumeItemListStart
            \resumeItem{Built ML data pipelines and statistical dependency estimates for stock pair selection; developed and backtested PyTorch trading strategies, outperforming baseline by 40\%.}
          \resumeItemListEnd

    \resumeSubHeadingListEnd

\section{Technical Skills}
 \begin{itemize}[leftmargin=0.15in, label={}]
    \small{\item{
     \textbf{Languages}{: Python, C/C++, Rust, Go, \LaTeX} \\
     \textbf{Tools}{: Linux, Docker, Kafka, Git, GitHub Actions, CMake, Yocto, BusyBox} \\
     \textbf{Data / ML}{: CUDA, NumPy, pandas, PyTorch, scikit-learn, SQL, DuckDB, Apache Arrow, Matplotlib/Seaborn}
    }}
 \end{itemize}

\end{document}
